\begin{helper}[The second recoloring has the full replacement suffix]
Fix a defective color $\delta$ and a
\noderef{ScheideDefectiveReturnData} $R$.  Write
$\gamma=\gamma_{\gammaIndex(\delta)}$, $\eta=\gamma_{R.j}$, and set
\[
 \psi_1=\noderef{OutsideTreeColorSwap}(T,\varphi,\gamma,\eta),
 \qquad
 K=\noderef{TwoColorComponent}
      (G,\psi_1,\eta,\delta,\operatorname{inl}R.v_1),
\]
\[
 \psi_2=\noderef{TwoColorComponentSwap}
      (G,\psi_1,\eta,\delta,\operatorname{inl}R.v_1).
\]
Then $\psi_1$ is proper, leaves exactly $e_0$ uncolored, is
$(T,\varnothing,\varphi)$-stable, and fixes every edge identity incident
with $T$.  The set $K$ is
exactly the set of nodes of the recorded replacement path at or after
the stored position of $R.v_1$; it contains $R.f_1$, meets $V(T)$ only
at $R.v_1$, and contains no displayed edge identity of $T$.  Moreover
$\psi_1(R.f_1)=\delta$.

The coloring $\psi_2$ is proper and leaves exactly $e_0$ uncolored;
$\delta$ is missing at $R.v_1$, every other vertex of $V(T)$ has its
original missing-color set, and
$\psi_2(f^0_\delta)=\delta$.
\end{helper}
\begin{proof}
Put $X=\noderef{TreeSequenceVertexSet}(T)$, let $P'$ be the full
replacement path recorded by \noderef{ScheideDefectiveSetup}, let $r$ be
the position of $R.v_1$, and let $S$ be the set of all vertex- and
edge-nodes of $P'$ at positions at least $r$.  By
\noderef{ScheideDefectiveReturnData}, the node at $r$ is $R.v_1$, the
next two nodes are $R.f_1,R.u_1$, and every vertex-node of $S$ other
than $R.v_1$ lies outside $X$.  The same datum gives
$R.v_1\in X$, $R.v_1\ne v^0_\delta$,
$\varphi(R.f_1)=\delta$, and
$\operatorname{ends}(R.f_1)=\{R.v_1,R.u_1\}$.  It also gives the three
pairwise inequalities among $\gamma,\eta,\delta$ and says that
$\eta$ is missing at $R.v_1$.

The \noderef{ScheideStructuralConfiguration} contained in the reduced
configuration contains a \noderef{GlobalMaximumTriple}.  Expanding its
\noderef{KTriple}, the coloring $\varphi$ is proper in the sense of
\noderef{IsProperEdgeColoring} and leaves exactly $e_0$ uncolored in the
sense of \noderef{IsOneEdgeUncolored}.  By
\noderef{ScheideGlobalMaximumOrder}, $X$ is
\noderef{ColoredVertexSetClosed}.  Moreover, $\gamma$ is missing at
$v^0_\delta$ and $\eta$ is missing at $R.v_1$; hence both belong to
\noderef{TreeMissingColors}$(\varphi,T)$.  The stability lemma
\noderef{OutsideTreeColorSwapStable}, applied to the distinct colors
$\gamma,\eta$ and the empty protected set, now proves that $\psi_1$ is
proper, leaves exactly $e_0$ uncolored, is
$(T,\varnothing,\varphi)$-\noderef{StableColoring}, and fixes every exact
edge identity incident with $T$.  This agrees directly with the defining
rule of \noderef{OutsideTreeColorSwap}.  In particular $R.f_1$ is fixed,
because it is incident with $R.v_1\in X$, so
$\psi_1(R.f_1)=\delta$.

We next prove the exact equality $S=K$.  Under $\varphi$, the support of
$P'$ is the full $(\gamma,\delta)$-component through
$v^0_\delta$ by \noderef{ScheideDefectiveChainPath}; in particular $P'$
is simple.  We first record that two nodes of $P'$ adjacent in the old
two-color incidence graph occur at consecutive positions.  At position
zero the vertex $v^0_\delta$ misses $\gamma$.  Thus every old two-color
neighbor there is a $\delta$-edge-node, and properness permits at most one
such exact identity: two distinct ones would be
\noderef{AdjacentEdges}, witnessed by their common
\noderef{Incident} vertex, and would have the same color.  At every later
position $n$ which has a later neighbor, the preceding and following
path nodes already give two distinct neighbors.  A third, nonconsecutive
neighbor would contradict
\noderef{TwoColorIncidenceDegreeAtMostTwo}; simplicity of $P'$ ensures
that the three nodes are distinct.  Reversing the order gives the claim
in both directions.

Every consecutive pair in the suffix $S$ is adjacent in the new
$(\eta,\delta)$ incidence graph.  The first edge is $R.f_1$ and was just
shown to retain color $\delta$.  At every later vertex, which lies outside
$X$, an incident old $\gamma$- or $\delta$-edge is handled as follows.
If its identity is not incident with $T$, the definition of the outside
swap changes $\gamma$ to $\eta$ and fixes $\delta$.  If it is incident
with $T$, choose its incident tree vertex.  Together with the displayed
outside endpoint, the endpoint-pair rule in
\noderef{LooplessMultigraph} and the definition of
\noderef{Incident} put this exact identity in
\noderef{BoundaryEdges}$(G,X)$.  Since $\gamma$ is missing at
$v^0_\delta\in X$, closedness rules out old color $\gamma$ on that
boundary identity; it therefore has color $\delta$ and is fixed.  This
proves consecutively that every node of $S$ is connected to $R.v_1$ in
the new incidence graph, so $S\subseteq K$ by the definition of
\noderef{TwoColorComponent}.

For completeness, the defective-color hypothesis says, through
\noderef{DefectiveColors} and \noderef{BoundaryColorEdges}, that there
are at least two $\delta$-colored boundary identities.  Hence this set is
nonempty; closedness then also shows that $\delta$ is present at every
vertex of $X$, since a vertex of $X$ missing $\delta$ would contradict
any one of those boundary identities.  This is consistent with, and
independently checks, the unique $\delta$-edge assertions used above.

To prove the reverse inclusion, we show that $S$ is closed under every
incidence adjacency in the new two-color graph.  Let the node at position
$n\ge r$ be newly adjacent to a node $z$.  If the first node is a
vertex-node, inspect the incident edge identity.  At $R.v_1$, an
$\eta$-edge is impossible because $\eta$ is missing there; at a later
outside vertex, a tree-incident $\eta$-edge would be a boundary identity
and would contradict closedness at $R.v_1$.  On a non-tree-incident edge,
the outside swap is an involution on $\gamma,\eta$.  Thus in all cases a
new $\eta$- or $\delta$-adjacency gives an old
$\gamma$- or $\delta$-adjacency.  If the first node is an edge-node, use
its preceding path vertex: its old affected color is read from that old
path adjacency, and the incidence of $z$ with the same exact edge gives
the required old adjacency.  Hence $z$ lies on $P'$, say at position
$m$, and the preceding consecutive-position claim gives
$m=n-1$ or $m=n+1$.

Both positions remain at least $r$ except possibly $n=r$ and $m=r-1$.
In that exceptional case the node at $n$ is $R.v_1$.  The predecessor
edge is incident with $X$ and is therefore fixed.  New adjacency says its
fixed color is $\eta$ or $\delta$; missingness of $\eta$ at $R.v_1$
forces color $\delta$.  It is distinct from the following edge $R.f_1$
by simplicity of $P'$, but both are incident with $R.v_1$ and both have
color $\delta$, contradicting properness.  Thus the exceptional adjacency
does not occur, and every new adjacency out of $S$ stays in $S$.  Induction
along any incidence walk from $R.v_1$ now gives $K\subseteq S$.  Therefore
$S=K$.

The stored position equalities put the edge-node $R.f_1$ in $S=K$.
The last-return condition gives directly
$K\cap X=\{R.v_1\}$ on vertex-nodes.  We prove separately, as required,
that $K$ contains no displayed edge identity of $T$.  Suppose the
edge-node of a displayed edge $e$ belonged to $K$.  Since the root of
$K$ is the distinct vertex-node $R.v_1$, component membership supplies a
nontrivial incidence path to that edge-node.  Hence the edge-node is in
the support of \noderef{TwoColorIncidenceGraph} and
$\psi_1(e)\in\{\eta,\delta\}$; this step excludes the possibility of an
isolated off-color edge-node.  By \noderef{TreeSequence}, the displayed
edge has as endpoints its newly displayed tree vertex and an earlier tree
vertex.  Each endpoint is incident with $e$, so the adjacency rule of the
two-color incidence graph puts both endpoint vertex-nodes in $K$.  The
singleton intersection just proved identifies both endpoints with
$R.v_1$.  Thus $e$ has endpoint pair
$\{R.v_1,R.v_1\}$, contradicting the loopless law in
\noderef{LooplessMultigraph}.  This proves all assertions about $K$.

The colors $\eta$ and $\delta$ are distinct, so
\noderef{TwoColorComponentSwapProper} makes $\psi_2$ proper.  The
definition of \noderef{TwoColorComponentSwap} merely exchanges two
nonempty color values on selected exact identities and fixes all other
values; it never changes whether a value is uncolored.  The two clauses of
\noderef{IsOneEdgeUncolored} therefore show that $\psi_2$ also leaves
exactly $e_0$ uncolored.  Stability of $\psi_1$ says that $\eta$ is still
missing at $R.v_1$.  Since that vertex-node belongs to $K$,
\noderef{TwoColorComponentSwapMissingBehavior} makes $\delta$ missing
there after the interchange.  If $v\in X$ and $v\ne R.v_1$, then the
singleton intersection gives $v\notin K$; the same missing-behavior lemma
first identifies its $\psi_2$- and $\psi_1$-missing sets, and stability
then identifies the latter with its original $\varphi$-missing set.

Finally $f^0_\delta$ is incident with
$v^0_\delta\in X$, so the first outside swap fixes it, and
\noderef{ScheideDefectiveSetup} gives
$\psi_1(f^0_\delta)=\delta$.  If its edge-node belonged to $K$, its
$\delta$ color and its incidence with $v^0_\delta$ would also put the
vertex-node $v^0_\delta$ in $K$.  The singleton intersection would give
$v^0_\delta=R.v_1$, contrary to the return datum.  Hence the edge-node of
$f^0_\delta$ is outside $K$, the second swap fixes this exact identity,
and $\psi_2(f^0_\delta)=\delta$.
\end{proof}
